% TOP_PROBLEM: {"id": 3000026, "title": "Deciding the validity of the score sequence of a soccer tournament", "category_id": 2, "category": {"id": 2, "name": "combinatorics", "display_name": "Combinatorics", "description": "Counting problems, graph theory, discrete structures.", "slug": "combinatorics", "order_index": 2, "created_at": "2026-07-31T15:26:25.671Z"}, "status": "partially_solved", "problem_number": "AMR-029-0026", "set_id": 13}
% FIRST_POSTED: 2026-10-01
% ENTRY_KIND: statement_only
% UnsolvedMath ID 3000026; source catalogue code AMR-029-0026
% !TeX program = lualatex
% Definition and sources only; this file is not an LLM solution attempt.
\documentclass[11pt,a4paper]{article}
\usepackage[margin=25mm]{geometry}
\usepackage{fontspec}
\setmainfont{lmroman10-regular.otf}[BoldFont=lmroman10-bold.otf,
 ItalicFont=lmroman10-italic.otf,BoldItalicFont=lmroman10-bolditalic.otf]
\usepackage{amsmath,amssymb,mathtools}
\usepackage{unicode-math}
\setmathfont{latinmodern-math.otf}
\AtBeginDocument{\let\setminus\smallsetminus}
\usepackage{xurl,hyperref}
\hypersetup{colorlinks=true,urlcolor=blue}
\setcounter{secnumdepth}{0}
\setlength{\parindent}{0pt}
\setlength{\parskip}{5pt}
\setlength{\emergencystretch}{3em}
\begin{document}
\section*{Deciding the validity of the score sequence of a soccer tournament}\label{3000026}
{\small UnsolvedMath ID 3000026 · AMR-029-0026 · Combinatorics · AMR Open Problem Lists}
\subsection{Definitions and mathematical statement}
Fix $n\ge1$ labelled teams. Each unordered pair plays exactly once.
The possible contributions of the match between $i$ and $j$ to their
respective scores are $(3,0)$, $(0,3)$, or $(1,1)$, corresponding to
a win, a loss, or a draw. Given nonnegative integers $s_1,\ldots,s_n$,
decide whether all matches can be assigned outcomes so that team $i$
receives exactly $s_i$ points.

More formally, for each $i<j$ choose
$(a_{ij},a_{ji})\in\{(3,0),(0,3),(1,1)\}$ and require
\[
                             \sum_{j\ne i}a_{ij}=s_i
                                  \quad(1\le i\le n).
\]
Is there a deterministic polynomial-time algorithm for this decision
problem, and for producing an outcome table when one exists?
The scores are given in binary; immediate infeasibility bounds such as
$s_i\le3(n-1)$ can be checked first.

No match outcomes are fixed beforehand. Allowing an arbitrary partially
played schedule as additional input is a different, more constrained
problem. The total score determines the number of drawn matches:
if $m=\binom n2$, that number must equal $3m-\sum_i s_i$.
This necessary aggregate condition does not specify how the individual
scores can be simultaneously realized by pairwise matches.
\subsection{Short English statement}
Can one efficiently decide whether prescribed team scores are
possible in a complete soccer tournament with three points for a win and one for a draw?
\subsection{Sources}
\begin{itemize}
\item[S1] ULAM.AI and the UnsolvedMath Contributors, supplied problems.json, record 3000026 (AMR-029-0026), AMR Open Problem Lists. Catalogue identity and source transcription; CC BY 4.0.
\url{https://huggingface.co/datasets/ulamai/UnsolvedMath}.
\item[S2] Egres Open - List of open problems, Open-problem page: Deciding the validity of the score sequence of a soccer tournament. Original source indexed by AMR-029-0026.
\url{https://oldlemon.cs.elte.hu/egres/open/List_of_open_problems}.
\item[S3] EGRES Open, Deciding the validity of the score sequence of a soccer tournament. Individual source page and explanatory remarks.
\url{https://lemon.cs.elte.hu/egres/open/Deciding_the_validity_of_the_score_sequence_of_a_soccer_tournament}.
\end{itemize}
\end{document}
